pdfoutput=1
% Options for packages loaded elsewhere
\PassOptionsToPackage{unicode}{hyperref}
\PassOptionsToPackage{hyphens}{url}
\pdfoutput=1
\documentclass[
  12pt,
]{article}
\usepackage{xcolor}
\usepackage{amsmath,amssymb}
\setcounter{secnumdepth}{-\maxdimen} % remove section numbering
\usepackage{iftex}
\ifPDFTeX
  \usepackage[T1]{fontenc}
  \usepackage{textcomp} % provide euro and other symbols
\else % if luatex or xetex
  \usepackage{unicode-math} % this also loads fontspec
  \defaultfontfeatures{Scale=MatchLowercase}
  \defaultfontfeatures[\rmfamily]{Ligatures=TeX,Scale=1}
\fi
\usepackage{lmodern}
\ifPDFTeX\else
  % xetex/luatex font selection
\fi
% Use upquote if available, for straight quotes in verbatim environments
\IfFileExists{upquote.sty}{\usepackage{upquote}}{}
\IfFileExists{microtype.sty}{% use microtype if available
  \usepackage[]{microtype}
  \UseMicrotypeSet[protrusion]{basicmath} % disable protrusion for tt fonts
}{}
\makeatletter
\@ifundefined{KOMAClassName}{% if non-KOMA class
  \IfFileExists{parskip.sty}{%
    \usepackage{parskip}
  }{% else
    \setlength{\parindent}{0pt}
    \setlength{\parskip}{6pt plus 2pt minus 1pt}}
}{% if KOMA class
  \KOMAoptions{parskip=half}}
\makeatother
\usepackage{longtable,booktabs,array}
\usepackage{calc} % for calculating minipage widths
% Correct order of tables after \paragraph or \subparagraph
\usepackage{etoolbox}
\makeatletter
\patchcmd\longtable{\par}{\if@noskipsec\mbox{}\fi\par}{}{}
\makeatother
% Allow footnotes in longtable head/foot
\IfFileExists{footnotehyper.sty}{\usepackage{footnotehyper}}{\usepackage{footnote}}
\makesavenoteenv{longtable}
\setlength{\emergencystretch}{3em} % prevent overfull lines
\providecommand{\tightlist}{%
  \setlength{\itemsep}{0pt}\setlength{\parskip}{0pt}}
\usepackage{bookmark}
\IfFileExists{xurl.sty}{\usepackage{xurl}}{} % add URL line breaks if available
\urlstyle{same}
\hypersetup{
  hidelinks,
  pdfcreator={LaTeX via pandoc}}

\author{SCX}
\date{2026}

\begin{document}

\section{Rigorous Mathematical Analysis of Multi-Head Spring and Physical Positional Encoding}\label{multi-head-spring-and-ppe-analysis}

\subsection{Honest Review of SCX Framework's Theorem Foundation}\label{honest-review-scx-framework}

\begin{center}\rule{0.5\linewidth}{0.5pt}\end{center}

\textbf{Document Status}: Complete analysis\\
\textbf{Dependencies}: Theorem 1--4 (theory/theorems/), SE-1, SE-2 (theory/self\_evolution/)\\
\textbf{Methodology}: Rigorous formal derivation + honest critique (no euphemisms)

\begin{center}\rule{0.5\linewidth}{0.5pt}\end{center}

\subsection{0. Notation Conventions and Framework Reconstruction}\label{notation-conventions}

Before entering the analysis, we precisely reconstruct the mathematical foundation of the SCX framework. All definitions below are rigorously formalized from the theorem descriptions in the prompt.

\subsubsection{0.1 State Crystallization}\label{state-crystallization}

Let the raw physical observation space be $\mathcal{X} \subset \mathbb{R}^{d_{\text{raw}}}$ (continuous physical quantities: PBE energy surface, force field, electron density, etc.). State Crystallization is a mapping:
\[\mathcal{C}: \mathcal{X} \to \mathcal{S} = \{s_1, \ldots, s_N\}\]
where each $s_i \in \mathbb{R}^{d_s}$ is a vector representation of a \textbf{discrete state atom}. Unlike BPE's statistical driving, $\mathcal{C}$ is physically driven: it clusters according to natural curvature boundaries of the PBE energy surface. Formally:
\[s_i = \mathcal{C}(x) \quad \text{where} \quad x \in \mathcal{X} \text{ and } \nabla^2 E_{\text{PBE}}(x) \text{ has a spectral gap at } x\]

\subsubsection{0.2 Yajie Audit}\label{yajie-audit}

Given $M$ independent expert models $\{E_1, \ldots, E_M\}$. For state atom $s_i$ and its label $y_i$, define:
\[c_m(s_i) = \mathbf{1}[E_m(s_i) \neq y_i] \quad \text{(expert $m$ disagrees with the given label)}\]

Consistency count:
\[C(s_i) = \sum_{m=1}^{M} c_m(s_i) \sim \text{Binomial}(M, p_{s_i})\]
where $p_{s_i}$ is the probability that experts ``disagree'' with the label. For clean samples, $p_{s_i} = p_{\text{clean}} < 0.5$; for noisy samples, $p_{s_i} = p_{\text{noisy}} > 0.5$.

\textbf{Cercis Score}: \[S(s_i) = Q(s_i) + \eta(t) \cdot N(s_i)\]
where $Q(s_i)$ is the base quality score (based on multi-expert consensus), $N(s_i)$ is the novelty bonus, and $\eta(t)$ is the time-decaying exploration weight.

\subsubsection{0.3 Spring Self-Evolution Dynamics}\label{spring-self-evolution-dynamics}

State transition: \[(S_t, \theta_t, M_t) \to (S_{t+1}, \theta_{t+1}, M_{t+1})\]

Spring self-attention form (single head):
\[\text{Spring}(s_i, \mathcal{S}) = \text{softmax}\left(\frac{(Q s_i)(K \mathcal{S})^T}{\sqrt{d_k}}\right) \cdot (V \mathcal{S})\]
where $Q, K \in \mathbb{R}^{d_k \times d_s}$, $V \in \mathbb{R}^{d_v \times d_s}$.

\begin{center}\rule{0.5\linewidth}{0.5pt}\end{center}

\subsection{1. Precise Theorem Statements of Core Theorems}\label{precise-core-theorem-statements}

\subsubsection{Theorem 1 (Multi-Expert Consensus Noise Detection)}\label{theorem-1-multi-expert}

\[\boxed{F_1 \geq 1 - \frac{1}{\eta} \sum_{s \in \mathcal{S}} \rho_s \cdot \exp\left(-2M \Delta_s^2\right)}\]
where:
- $\rho_s = |\{i: \mathcal{C}(x_i) = s\}| / N_{\text{total}}$: fraction of state atom $s$.
- $\Delta_s = |p_{\text{clean},s} - p_{\text{noisy},s}|$: detection margin for state $s$.
- $M$: number of independent experts.
- $\eta > 0$: normalization factor.

\textbf{Proof technique}: Chernoff bound + Hoeffding inequality. For each state $s$, $C(s)$ is a mixture of two binomial distributions. Mark as noise when $C(s) \geq \tau_s$. Chernoff yields:
\[P(\text{False Positive}) = P_{H_0}(C(s) \geq \tau_s) \leq \exp(-2M(\tau_s/M - p_{\text{clean},s})^2)\]
\[P(\text{False Negative}) = P_{H_1}(C(s) < \tau_s) \leq \exp(-2M(p_{\text{noisy},s} - \tau_s/M)^2)\]
Choose $\tau_s = M \cdot (p_{\text{clean},s} + p_{\text{noisy},s})/2$ balancing both terms, obtaining $\Delta_s$.

\subsubsection{Theorem 2 (Inevitable Failure of Weak Features)}\label{theorem-2-weak-features}

\[\boxed{F_{1,\text{SCX}} \leq F_{1,\text{base}} + C_F \cdot \sqrt{\frac{\delta}{2}}}\]
where $\delta > 0$ characterizes the information deficiency of the feature space.

\textbf{Proof technique}: Fano's inequality. Define $\delta = \min_{f} \mathbb{E}[\mathbf{1}[f(X) \neq Y]] - \text{Bayes error}$, the minimum gap between the best feasible classifier and the Bayes optimal. Fano relates conditional entropy to error rate:
\[P_e \geq \frac{H(Y|X) - 1}{\log |\mathcal{Y}|}\]
When $H(Y|X)$ is large (weak features), any method's (including SCX's) error rate lower bound is determined by $\delta$.

\subsubsection{Theorem 3 (Indistinguishability of Noise and Difficult Samples)}\label{theorem-3-indistinguishability}

\textbf{Constructive two-world proof}: There exist two worlds $W_A, W_B$ with observationally equivalent data distributions $P_{W_A}(X, Y) = P_{W_B}(X, Y)$, but:
\begin{itemize}
\tightlist
\item
  World $W_A$: sample $(x, y)$ is mislabeled (label noise).
\item
  World $W_B$: sample $(x, y)$ is correctly labeled but intrinsically difficult (high Bayesian uncertainty).
\end{itemize}
No algorithm can distinguish these two worlds from observational data alone.

Formally: $\forall$ decision rule $d: (\mathcal{X} \times \mathcal{Y})^n \to \{W_A, W_B\}$,
\[\max_{W \in \{W_A, W_B\}} P_W(d(\text{data}) \neq W) \geq \frac{1}{2} - \frac{1}{2}\|P_{W_A}^{\otimes n} - P_{W_B}^{\otimes n}\|_{\text{TV}} = \frac{1}{2}\]
since $P_{W_A}^{\otimes n} = P_{W_B}^{\otimes n}$ (completely identical observation distributions).

\subsubsection{Theorem 4 (Exact Constant Minimax Optimality)}\label{theorem-4-exact-constant-minimax}

SCX's noise detection error rate achieves exact constant optimality in the following sense:
\[\lim_{M \to \infty} -\frac{1}{M} \log P_{\text{FN}}^{(M)} = D_{\text{KL}}(P_{\text{noisy}} \| P_{\text{clean}}) \quad \text{subject to } P_{\text{FP}}^{(M)} \leq \alpha\]
and the $O(1/\sqrt{M})$ correction term from the Bahadur-Rao expansion also matches the minimax lower bound.

\textbf{Proof technique}: Chernoff-Stein lemma + Bahadur-Rao expansion.

\begin{center}\rule{0.5\linewidth}{0.5pt}\end{center}

\subsection{2. Component 1: Mathematical Analysis of Physical Positional Encoding (PPE)}\label{component-1-ppe-analysis}

\subsubsection{2.1 Formal Definition}\label{ppe-formal-definition}

Let each state atom $s_i$ carry a physical position $p_i \in \mathcal{P}$, where:
- Protein sequence: $\mathcal{P} = \mathbb{N}$ (residue index).
- 3D structure: $\mathcal{P} = \mathbb{R}^3$ (atom coordinates).
- Docking: $\mathcal{P} = \mathbb{R}^3$ (binding pocket coordinates).

\textbf{Definition 1 (Physical Positional Encoding)}: PPE is a mapping
\[\text{PE}: \mathcal{P} \to \mathbb{R}^{d_{\text{pe}}}\]
The encoded state representation: \[\boxed{h_i = \phi(s_i) + \text{PE}(p_i)}\]
where $\phi: \mathbb{R}^{d_s} \to \mathbb{R}^{d}$ is a feature map, and we require $d_{\text{pe}} = d$.

\textbf{Alternative (concatenation + projection)}:
\[h_i = W \cdot [\phi(s_i); \text{PE}(p_i)]\]
where $W \in \mathbb{R}^{d \times (d + d_{\text{pe}})}$. Mathematically this is more general, but the additive scheme is cleaner and degenerates to addition when $W = [I_d, I_d]$. The analysis below focuses on the additive scheme; results can be generalized.

\subsubsection{2.2 Concrete Forms of Encoding Functions}\label{encoding-function-forms}

\textbf{Sine encoding for scalar positions} (analogous to Transformer):
\[\text{PE}(p, 2j) = \sin\left(\frac{p}{10000^{2j/d}}\right), \quad \text{PE}(p, 2j+1) = \cos\left(\frac{p}{10000^{2j/d}}\right)\]
where $j = 0, 1, \ldots, d/2 - 1$.

\textbf{Rotation encoding for 3D positions} (analogous to RoPE, adapted to physical space):
\[\text{PE}_{\text{rot}}(\mathbf{p}) = \mathbf{R}(\mathbf{p}) \cdot \mathbf{e}_0\]
where $\mathbf{R}(\mathbf{p}) \in SO(d)$ is a rotation matrix parameterized by coordinates. For $\mathbf{p} = (x, y, z)$:
\[\mathbf{R}(\mathbf{p}) = \mathbf{R}_x(\omega_x x) \cdot \mathbf{R}_y(\omega_y y) \cdot \mathbf{R}_z(\omega_z z)\]
where $\mathbf{R}_\alpha(\theta)$ is a $d$-dimensional rotation in the $\alpha$-axis plane, and $\omega_x, \omega_y, \omega_z$ are learnable or fixed frequency parameters.

\textbf{Key property}: Unlike NLP positional encoding, physical positions possess \textbf{multi-dimensional geometric structure}. A good PPE should respect physical symmetries:

\begin{enumerate}
\def\labelenumi{\arabic{enumi}.}
\tightlist
\item
  \textbf{Translation equivariance} (for 3D structures): $\text{PE}(\mathbf{p} + \mathbf{t})$ has a predictable relationship with $\text{PE}(\mathbf{p})$.
\item
  \textbf{Rotation equivariance}: $\text{PE}(\mathbf{R}\mathbf{p})$ can be recovered by rotating the encoding space.
\item
  \textbf{Lipschitz continuity}: $\|\text{PE}(\mathbf{p}) - \text{PE}(\mathbf{q})\| \leq L_{\text{PE}} \cdot \|\mathbf{p} - \mathbf{q}\|$.
\end{enumerate}

\subsubsection{2.3 Impact on Theorem 1}\label{impact-on-theorem-1}

\textbf{Core question}: How does PPE change the detection margin $\Delta_s$?

\textbf{Proposition 2.1 (Modified Detection Margin under PPE)}: In the PPE-augmented feature space, the effective detection margin for state $s$ is:
\[\boxed{\Delta_s^{\text{PPE}} = \Delta_s + \delta_s^{\text{PE}}}\]
where $\delta_s^{\text{PE}} \in [-\Delta_s, 1 - p_{\text{clean},s} - \Delta_s]$ satisfies:
\[\delta_s^{\text{PE}} = \underbrace{(p_{\text{clean},s}^{\text{PPE}} - p_{\text{clean},s})}_{\text{consistency gain on clean samples}} - \underbrace{(p_{\text{noisy},s}^{\text{PPE}} - p_{\text{noisy},s})}_{\text{confusion gain on noisy samples}}\]

\emph{Proof}: In the augmented space, expert prediction probabilities become $p_{\text{clean},s}^{\text{PPE}} = \mathbb{E}[\mathbf{1}[E_m(h_i) = y_i] | H_0]$ and $p_{\text{noisy},s}^{\text{PPE}} = \mathbb{E}[\mathbf{1}[E_m(h_i) \neq y_i] | H_1]$. The detection margin is defined as $\Delta_s^{\text{PPE}} = |p_{\text{clean},s}^{\text{PPE}} - (1 - p_{\text{noisy},s}^{\text{PPE}})| = |p_{\text{clean},s}^{\text{PPE}} + p_{\text{noisy},s}^{\text{PPE}} - 1|$. Expansion yields the result. $\square$

\textbf{Theorem 2.1 (F1 Lower Bound under PPE)}:
\[\boxed{F_1^{\text{PPE}} \geq 1 - \frac{1}{\eta} \sum_{s \in \mathcal{S}} \rho_s \cdot \exp\left(-2M (\Delta_s + \delta_s^{\text{PE}})^2\right)}\]
The \textbf{validity of the Chernoff-Hoeffding inequality is unchanged} (indicator variables remain bounded in $[0,1]$), but the constant changes.

When $\delta_s^{\text{PE}} > 0$ for all $s$ (PPE always helps detection):
\[F_1^{\text{PPE}} - F_1 \geq \frac{1}{\eta} \sum_s \rho_s \left[\exp(-2M\Delta_s^2) - \exp(-2M(\Delta_s + \delta_s^{\text{PE}})^2)\right] \geq 0\]

\textbf{Proposition 2.2 (Spatial Locality Regularization)}: If PE satisfies the Lipschitz condition (Lipschitz constant $L_{\text{PE}}$) and expert model $E_m$ is $L_E$-Lipschitz in its input, then for any two state atoms with physical positions $p_i, p_j$:
\[|\delta_i^{\text{PE}} - \delta_j^{\text{PE}}| \leq 2L_E \cdot L_{\text{PE}} \cdot \|p_i - p_j\|\]

\emph{Proof}: Changes in expert consistency probability are constrained by input changes:
\begin{align*}
|p_{\text{clean},i}^{\text{PPE}} - p_{\text{clean},j}^{\text{PPE}}|
&\leq \mathbb{E}_{E} \left|\mathbf{1}[E(h_i) = y] - \mathbf{1}[E(h_j) = y]\right| \\
&\leq \mathbb{E}_{E} \left[ L_E \cdot \|h_i - h_j\| \right] \\
&= L_E \cdot \|\text{PE}(p_i) - \text{PE}(p_j)\| \quad (\text{since } s_i = s_j = s) \\
&\leq L_E \cdot L_{\text{PE}} \cdot \|p_i - p_j\|
\end{align*}
The same holds for $p_{\text{noisy}}$. The difference of the two yields the factor 2. $\square$

This means \textbf{physically neighboring state atoms have similar detection difficulty} --- a verifiable prediction.

\subsubsection{2.4 Impact on Theorem 2}\label{impact-on-theorem-2}

\textbf{Core question}: How does imperfect PPE change the weak feature bound?

\textbf{Definition 2 (Encoding Imperfection $\varepsilon_{\text{PE}}$)}:
\[\boxed{\varepsilon_{\text{PE}} = I(Y; P | X) - I(Y; \text{PE}(P) | X) \geq 0}\]
where $I(Y; P | X)$ is the conditional mutual information between physical position $P$ and label $Y$ given raw features $X$, and $I(Y; \text{PE}(P) | X)$ is the information retained after encoding. $\varepsilon_{\text{PE}}$ quantifies the \textbf{loss of positional information during encoding}.

\textbf{Theorem 2.2 (Failure Upper Bound under Imperfect PPE)}:
\[\boxed{F_{1,\text{SCX}}^{\text{PPE}} \leq F_{1,\text{base}} + C_F \cdot \sqrt{\frac{\delta + \frac{2\varepsilon_{\text{PE}}}{C_F^2}}{2}}}\]

\emph{Proof}: Fano's inequality yields $P_e^{\text{PPE}} \geq \frac{H(Y|X, \text{PE}(P)) - 1}{\log |\mathcal{Y}|}$. We have:
\begin{align*}
H(Y|X, \text{PE}(P)) &= H(Y|X) - I(Y; \text{PE}(P) | X) \\
&= H(Y|X) - [I(Y; P | X) - \varepsilon_{\text{PE}}] \\
&\geq H(Y|X) - I(Y; P | X) + \varepsilon_{\text{PE}}
\end{align*}
The definition of $\delta$ yields $H(Y|X) - 1 \geq (\delta/2) \cdot \log|\mathcal{Y}|$ (via the inverse direction of Fano). Therefore:
\[H(Y|X, \text{PE}(P)) - 1 \geq \left(\frac{\delta}{2} + \frac{\varepsilon_{\text{PE}}}{C_F}\right) \cdot \log|\mathcal{Y}|\]
Substituting into SCX's F1 upper bound expression yields the result. $\square$

\textbf{Tightness analysis}:
- When $\varepsilon_{\text{PE}} = 0$ (perfect encoding): The upper bound recovers the original Theorem 2.
- When $\varepsilon_{\text{PE}} > 0$: The upper bound relaxes at rate $\sqrt{\delta + O(\varepsilon_{\text{PE}})}$.
- When $\varepsilon_{\text{PE}} \to I(Y; P | X)$ (completely failed encoding, lost all position information): The upper bound returns to the case of completely unused position information.

\textbf{Corollary 2.1 (Value of Position Information)}: If physical position $P$ itself carries no information about label $Y$ (i.e., $I(Y; P | X) = 0$), then $\varepsilon_{\text{PE}} = 0$ (since there is no information to lose), and Theorem 2's upper bound is \textbf{completely unchanged}. PPE in this case is \textbf{pure parameter waste}.

\subsubsection{2.5 Impact on Theorem 3}\label{impact-on-theorem-3}

This is the most critical analysis. Theorem 3's indistinguishability relies on constructing two worlds $W_A, W_B$ with completely identical observation distributions $P(X, Y)$.

\textbf{Proposition 2.3 (Fixed PPE Preserves Theorem 3)}: If PE is an \textbf{a priori fixed} function (independent of training data or labels), then the augmented feature distributions in the two worlds remain identical:
\[P_{W_A}(X, \text{PE}(P), Y) = P_{W_B}(X, \text{PE}(P), Y)\]
Therefore Theorem 3's indistinguishability \textbf{remains unchanged}. PPE neither helps nor hurts worst-case identifiability.

\emph{Proof}: Physical position $P$ is determined by physical structure, independent of the label generation process. Hence $P_{W_A}(P|X) = P_{W_B}(P|X)$. Since PE is a fixed deterministic function, $P_{W_A}(\text{PE}(P)|X) = P_{W_B}(\text{PE}(P)|X)$. The observation distributions are completely identical. $\square$

\textbf{Proposition 2.4 (Learned PPE Breaks Theorem 3)}: If PE's parameters $\theta_{\text{PE}}$ are learned from data (including labels $Y$), then different encodings are learned in the two worlds:
\[\text{PE}_{W_A} \neq \text{PE}_{W_B}\]
because $W_A$'s noisy labels affect the learning process. In this case:
\[P_{W_A}(X, \text{PE}_{W_A}(P), Y) \neq P_{W_B}(X, \text{PE}_{W_B}(P), Y)\]
The two worlds become \textbf{distinguishable}. Theorem 3 is broken.

\textbf{Honest assessment}: This is actually \textbf{good news}. Theorem 3 is an impossibility result --- it says that in the worst case, noise and difficult samples are indistinguishable. If learned PPE can break Theorem 3's premise, it means \textbf{position information provides the ability to distinguish noise from difficult samples}. This is not a flaw of the framework but progress of the framework. However, it does mean Theorem 3 no longer holds unconditionally --- one must add the premise ``PE is fixed or position information carries no label-relevant information.''

\subsubsection{2.6 Impact on Theorem 4}\label{impact-on-theorem-4}

\textbf{Proposition 2.5 (PPE Does Not Reduce the Optimal Error Exponent)}:
\[\boxed{D_{\text{KL}}(P_{\text{noisy}}^{\text{PPE}} \| P_{\text{clean}}^{\text{PPE}}) = D_{\text{KL}}(P_{\text{noisy}} \| P_{\text{clean}}) + \Delta D_{\text{KL}}^{\text{PE}}}\]
where $\Delta D_{\text{KL}}^{\text{PE}} \geq 0$ (non-negative, guaranteed by the data processing inequality for KL divergence).

\emph{Proof}: By the chain rule for KL divergence:
\begin{align*}
D_{\text{KL}}(P_{\text{noisy}}^{\text{PPE}} \| P_{\text{clean}}^{\text{PPE}})
&= D_{\text{KL}}(P_{\text{noisy}}(X, H) \| P_{\text{clean}}(X, H)) \\
&= D_{\text{KL}}(P_{\text{noisy}}(X) \| P_{\text{clean}}(X)) + \mathbb{E}_{X \sim P_{\text{noisy}}}[D_{\text{KL}}(P_{\text{noisy}}(H|X) \| P_{\text{clean}}(H|X))]
\end{align*}
Since $H = \phi(X) + \text{PE}(P)$ and PE is a fixed deterministic function, augmentation does not reduce KL divergence (information processing inequality). Specifically:
\[\Delta D_{\text{KL}}^{\text{PE}} = \mathbb{E}_{X \sim P_{\text{noisy}}}[D_{\text{KL}}(P_{\text{noisy}}(\text{PE}(P)|X) \| P_{\text{clean}}(\text{PE}(P)|X))] \geq 0\]
$\square$

PPE \textbf{never reduces} the optimal error exponent. If position information is label-relevant, $\Delta D_{\text{KL}}^{\text{PE}} > 0$, and the error exponent is \textbf{strictly larger} (more accurate detection). If position information is irrelevant, $\Delta D_{\text{KL}}^{\text{PE}} = 0$, unchanged.

The Bahadur-Rao constant is modified by the variance term introduced by PE, but always in the direction of tighter bounds (when PE is useful).

\begin{center}\rule{0.5\linewidth}{0.5pt}\end{center}

\subsection{3. Component 2: Mathematical Analysis of Multi-Head Spring}\label{component-2-multi-head-spring}

\subsubsection{3.1 Formal Definition}\label{multi-head-spring-definition}

\textbf{Definition 3 (Multi-Head Spring)}:
\[\boxed{\text{Spring}_{\text{MH}}(s_i, \mathcal{S}) = \text{Concat}(\text{head}_1, \ldots, \text{head}_K) \cdot W^O}\]
where:
\[\text{head}_k = \text{Spring}_k(s_i, \mathcal{S}) = \text{softmax}\left(\frac{(Q_k s_i)(K_k \mathcal{S})^T}{\sqrt{d_k}}\right) \cdot (V_k \mathcal{S})\]
Parameters:
- $Q_k, K_k \in \mathbb{R}^{d_k \times d_s}$: query/key projections for head $k$.
- $V_k \in \mathbb{R}^{d_v \times d_s}$: value projection for head $k$.
- $W^O \in \mathbb{R}^{K d_v \times d_s}$: output projection.
- $d_k = d_s / K$ (standard multi-head dimensionality reduction convention), or $d_k = d_s$ (full-dimensional multi-head, each head operates on the full state space).

Each head $k$ is designed to attend to a different physical dimension:
- Head 1: bond angle patterns --- attends to angular relationships among state atoms.
- Head 2: bond length patterns --- attends to distance relationships among state atoms.
- Head 3: coordination environment --- attends to local neighborhood topology of state atoms.
- Head 4: energy gradient / force field patterns --- attends to energy surface curvature of state atoms.

\textbf{Physical dimension constraint (structured multi-head)}: To achieve genuine specialization, each head's $Q_k, K_k$ should be constrained to act only on specific subspaces of the state vector. If $s_i = [s_i^{(1)}; s_i^{(2)}; s_i^{(3)}; s_i^{(4)}]$ is partitioned by physical dimension, then:
\[Q_k = [\mathbf{0}_{d_k \times d_{<k}} \;|\; Q_k^{(k)} \;|\; \mathbf{0}_{d_k \times d_{>k}}]\]
where $Q_k^{(k)}$ only acts on the $k$-th physical dimension subspace, with the rest zero. This \textbf{forces} specialization.

\subsubsection{3.2 Impact on Theorem 1 --- Effective Number of Experts $M$}\label{impact-on-thm-1-effective-M}

\textbf{Core question}: Are $K$ heads equivalent to $K$ independent experts?

\textbf{Answer: No.} This is a critical distinction.

\textbf{Proposition 3.1 (Heads Are Not Independent Experts)}: Multi-Head Spring's head outputs are \textbf{conditionally dependent} given the set of state atoms:
\[\text{Cov}(\text{head}_k(s_i), \text{head}_l(s_i) \mid \mathcal{S}) \neq 0 \quad \text{(in general)}\]

\emph{Proof}: All heads share the same input state atoms $\mathcal{S}$ and the same softmax attention pooling structure. Even if the projection matrices $Q_k, K_k$ differ, the attention weights $\alpha_{ij}^{(k)} = \text{softmax}_j((Q_k s_i)^T (K_k s_j) / \sqrt{d_k})$ are coupled through the common $s_j$ values. Therefore:
\[\text{head}_k(s_i) = \sum_j \alpha_{ij}^{(k)} \cdot V_k s_j\]
\[\text{head}_l(s_i) = \sum_j \alpha_{ij}^{(l)} \cdot V_l s_j\]
The covariance between the two is non-zero through the shared $\{s_j\}$ terms. $\square$

\textbf{Definition 4 (Effective Expert Diversity $\rho_{\text{eff}}$)}:
\[\boxed{\rho_{\text{eff}} = 1 - \frac{2}{K(K-1)} \sum_{k < l} |\text{Corr}(\text{head}_k, \text{head}_l)|}\]
where $\text{Corr}(\text{head}_k, \text{head}_l) = \frac{\text{Cov}(\text{head}_k, \text{head}_l)}{\sqrt{\text{Var}(\text{head}_k) \cdot \text{Var}(\text{head}_l)}}$ (element-wise average correlation).

\begin{itemize}
\tightlist
\item
  $\rho_{\text{eff}} = 1$: all heads completely independent (ideal).
\item
  $\rho_{\text{eff}} = 0$: all heads completely correlated (degenerates to single head).
\item
  $\rho_{\text{eff}} < 0$: negative correlation among heads (adversarial diversity).
\end{itemize}

\textbf{Heuristic correction for effective number of experts}:
\[M_{\text{eff}} = M \cdot (1 + \alpha \cdot \rho_{\text{eff}} \cdot (K - 1))\]
where $\alpha \in [0, 1]$ is the ``head diversity $\to$ expert diversity'' conversion coefficient.

\textbf{But this is not a rigorous bound}. A rigorous Chernoff bound requires handling dependent random variables.

\textbf{Theorem 3.1 (Rigorous F1 Lower Bound under Multi-Head Spring)}:
Under the dependence structure induced by shared state atoms, the F1 lower bound is:
\[\boxed{F_1^{\text{MH}} \geq 1 - \frac{1}{\eta} \sum_{s \in \mathcal{S}} \rho_s \cdot \exp\left(-2M_{\text{eff}} \Delta_s^2\right)}\]
where $M_{\text{eff}} = M / (1 + \delta_{\text{MH}})$ and $\delta_{\text{MH}} \geq 0$ depends on the strength of inter-head dependence. The exact form of $\delta_{\text{MH}}$ requires specifying the dependence structure (e.g., $\beta$-mixing coefficients, or using Janson's inequality for dependent indicators).

\emph{Proof sketch}: For dependent Bernoulli variables, Hoeffding's inequality is replaced by Janson's inequality or Azuma-Hoeffding for martingale difference sequences under $\beta$-mixing. The degradation factor $1/(1+\delta_{\text{MH}})$ captures the effective sample size reduction due to dependence. $\square$

\subsubsection{3.3 Impact on Spring Convergence (SE-1)}\label{impact-on-spring-convergence}

\textbf{Proposition 3.2 (Multi-Head Loss Landscape)}: The loss landscape of Multi-Head Spring possesses permutation symmetry: for any permutation $\pi \in S_K$ of head indices,
\[\mathcal{L}_{\text{Spring}}(\pi \cdot \theta) = \mathcal{L}_{\text{Spring}}(\theta)\]
This generates $K!$ equivalent stationary points (equivalence classes).

\textbf{Consequence}: The convergence analysis of the coupled $(S_t, \theta_t)$ system must account for permutation symmetry in the Lyapunov analysis. The Lyapunov function must be $S_K$-invariant to be well-defined on the quotient space $\Theta/S_K$.

\begin{center}\rule{0.5\linewidth}{0.5pt}\end{center}

\subsection{4. Cross-Component Interactions}\label{cross-component-interactions}

\subsubsection{4.1 PPE $\times$ Multi-Head Spring Combined F1 Bound}\label{ppe-x-multi-head-combined}

Combining Theorem 2.1 and Theorem 3.1:
\[\boxed{F_1^{\text{MH+PPE}} \geq 1 - \frac{1}{\eta} \sum_{s \in \mathcal{S}} \rho_s \cdot \exp\left(-2M_{\text{eff}} (\Delta_s + \delta_s^{\text{PE}})^2\right)}\]

\textbf{Interpretation}: PPE increases the per-state detection margin ($\Delta_s \to \Delta_s + \delta_s^{\text{PE}}$), while Multi-Head Spring adjusts the effective number of experts ($M \to M_{\text{eff}}$). The two effects are multiplicative in the exponent --- a positive $\delta_s^{\text{PE}}$ and a small $\delta_{\text{MH}}$ (near-independent heads) produce the strongest bound.

\subsubsection{4.2 Honest Summary}\label{honest-summary}

\begin{longtable}[]{@{}
  >{\raggedright\arraybackslash}p{(\linewidth - 4\tabcolsep) * \real{0.2500}}
  >{\raggedright\arraybackslash}p{(\linewidth - 4\tabcolsep) * \real{0.2500}}
  >{\raggedright\arraybackslash}p{(\linewidth - 4\tabcolsep) * \real{0.2500}}
  >{\raggedright\arraybackslash}p{(\linewidth - 4\tabcolsep) * \real{0.2500}}@{}}
\toprule\noalign{}
\begin{minipage}[b]{\linewidth}\raggedright
Component
\end{minipage} & \begin{minipage}[b]{\linewidth}\raggedright
Theoretical Benefit
\end{minipage} & \begin{minipage}[b]{\linewidth}\raggedright
Key Limitation
\end{minipage} & \begin{minipage}[b]{\linewidth}\raggedright
Status
\end{minipage} \\
\midrule\noalign{}
\endhead
\bottomrule\noalign{}
\endlastfoot
Fixed PPE & Increases $\Delta_s$ when $I(Y;P|S)>0$ & No benefit when $I(Y;P|S)=0$; rotation non-equivariance & \textbf{Conditionally proven} \\
Learned PPE & Can break Theorem 3 (good) & Requires auxiliary loss with physical quantities & \textbf{Open problem} \\
Multi-Head Spring & Increases effective model capacity & Heads are dependent (not independent experts) & \textbf{Partially rigorous} \\
PPE + MH combined & Multiplicative exponent improvement & Both limitations accumulate; rigorous bound needs dependence handling & \textbf{Heuristic bound} \\
\end{longtable}

\end{document}
